\documentclass[11pt,a4paper]{article}
\usepackage{shared/luxcover}
\usepackage[utf8]{inputenc}
\usepackage{amsmath,amssymb,amsfonts,amsthm}
\usepackage{graphicx}
\usepackage{hyperref}
\usepackage{algorithm}
\usepackage{algpseudocode}
\usepackage{xcolor}
\usepackage{booktabs}
\usepackage{enumitem}
\usepackage{tikz}

\hypersetup{colorlinks=true,linkcolor=black,citecolor=blue,urlcolor=blue}

\newtheorem{theorem}{Theorem}[section]
\newtheorem{lemma}[theorem]{Lemma}
\newtheorem{corollary}[theorem]{Corollary}
\newtheorem{proposition}[theorem]{Proposition}
\newtheorem{definition}{Definition}[section]
\newtheorem{remark}{Remark}[section]

\title{\textbf{Prism Protocol: Multi-Spectrum Parallel Consensus}}

\author{
Zach Kelling\thanks{Corresponding author: zach@lux.network}\\
\textit{Lux Industries}\\
\texttt{research@lux.network}
}

\date{2024}

\begin{document}

\luxcoverpage

\begin{abstract}
We present the Prism protocol, a multi-spectrum parallel consensus engine that
decomposes the consensus problem into independent frequency bands processed concurrently.
Prism assigns each transaction to a spectrum band based on its conflict set hash,
ensuring that conflicting transactions always land in the same band while independent
transactions are distributed across bands. Each band operates its own Photon pipeline
independently, and a \emph{spectral combiner} merges finalized results across bands
into a globally consistent total order.
We formalize the spectral decomposition using a hash-based partitioning scheme and
prove three key properties:
(1) \emph{Conflict locality}: all transactions in a conflict set map to the same band
with probability 1 (by construction).
(2) \emph{Load balance}: for uniformly distributed transaction hashes, the expected
load per band is $T/B$ where $T$ is the total transaction rate and $B$ is the number
of bands, with deviation $O(\sqrt{T \log B / B})$.
(3) \emph{Order consistency}: the spectral combiner produces a total order compatible
with the partial orders of all individual bands.
The protocol achieves aggregate throughput of $\Theta(B \cdot \Theta_{\text{band}})$
where $\Theta_{\text{band}}$ is the per-band Photon throughput, providing linear
scaling with the number of spectrum bands. With $B = 16$ bands and 1,000 validators,
Prism achieves 50,000 TPS with 500ms median finality per transaction.
Direct voting mode provides an $O(n)$ fast path for uncontested decisions.
\end{abstract}

\tableofcontents
\newpage

%% -----------------------------------------------------------------------
%%  1. INTRODUCTION
%% -----------------------------------------------------------------------
\section{Introduction}
\label{sec:introduction}

The throughput of a single consensus instance is fundamentally limited by the per-round
message complexity and the finality latency. Even with pipelining (Photon) and burst
finalization (Nova), a single consensus instance processes transactions sequentially
within each conflict set. For blockchains with diverse transaction types---token
transfers, smart contract calls, cross-chain messages, governance votes---the aggregate
throughput is limited by the busiest conflict set.

Prism addresses this by partitioning the consensus problem into independent spectrum
bands, each handling a subset of conflict sets. This is analogous to frequency-division
multiplexing in telecommunications: each band occupies a different ``frequency''
(conflict set partition) and processes transactions independently.

The challenge is ensuring global consistency: while each band finalizes independently,
the combined output must form a valid total order compatible with cross-band dependencies.
Prism's spectral combiner solves this through a deterministic merge algorithm that
respects the causal order of cross-band references.

\paragraph{Contributions.}
\begin{enumerate}[nosep]
\item The Prism multi-spectrum consensus protocol with hash-based band assignment (\S\ref{sec:protocol}).
\item Conflict locality and load balance proofs (\S\ref{sec:partition}).
\item The spectral combiner algorithm with order consistency proof (\S\ref{sec:combiner}).
\item Direct voting fast path for uncontested decisions (\S\ref{sec:voting}).
\item Linear throughput scaling demonstration: 50,000 TPS with 16 bands (\S\ref{sec:evaluation}).
\end{enumerate}

%% -----------------------------------------------------------------------
%%  2. SPECTRAL DECOMPOSITION
%% -----------------------------------------------------------------------
\section{Spectral Decomposition}
\label{sec:partition}

\subsection{Band Assignment}

\begin{definition}[Spectrum Band]
Given $B$ spectrum bands $\{b_0, b_1, \ldots, b_{B-1}\}$, a transaction $t$ with
conflict set identifier $\text{csid}(t)$ is assigned to band:
\[
\text{band}(t) = H(\text{csid}(t)) \mod B
\]
where $H$ is a cryptographic hash function (SHA-256).
\end{definition}

\begin{definition}[Conflict Set Identifier]
The conflict set identifier $\text{csid}(t)$ is determined by the consumed UTXO(s)
of transaction $t$. All transactions consuming the same UTXO share the same $\text{csid}$.
Non-conflicting transactions have distinct $\text{csid}$ values.
\end{definition}

\begin{theorem}[Conflict Locality]
\label{thm:locality}
If $t_1 \in \mathcal{C}(t_2)$ (the transactions conflict), then $\text{band}(t_1) = \text{band}(t_2)$.
Conflicting transactions are always processed in the same spectrum band.
\end{theorem}

\begin{proof}
By definition, $t_1 \in \mathcal{C}(t_2)$ implies $\text{csid}(t_1) = \text{csid}(t_2)$
(they consume the same UTXO). Therefore:
\[
\text{band}(t_1) = H(\text{csid}(t_1)) \mod B = H(\text{csid}(t_2)) \mod B = \text{band}(t_2)
\]
\end{proof}

\subsection{Load Balance}

\begin{theorem}[Load Balance]
\label{thm:balance}
For $T$ transactions with uniformly distributed $\text{csid}$ values and $B$ bands,
the load on band $b$ satisfies:
\[
\Pr\left[|L_b - T/B| > \sqrt{\frac{2T \ln B}{B}}\right] < \frac{2}{B}
\]
where $L_b$ is the number of transactions assigned to band $b$.
\end{theorem}

\begin{proof}
Each transaction is independently assigned to band $b$ with probability $1/B$ (by
the uniformity of $H$). The load $L_b \sim \text{Binomial}(T, 1/B)$ with mean $T/B$
and variance $T(B-1)/B^2$.

By the Chernoff bound:
\[
\Pr[|L_b - T/B| > \delta] \leq 2\exp\left(-\frac{\delta^2 B}{2T}\right)
\]
Setting $\delta = \sqrt{2T \ln B / B}$ gives $\Pr[|L_b - T/B| > \delta] < 2/B^2$.
A union bound over $B$ bands:
\[
\Pr[\exists b : |L_b - T/B| > \delta] < B \cdot 2/B^2 = 2/B
\]
For $B = 16$: the maximum deviation from mean load is at most $\sqrt{2T \ln 16 / 16} \approx 0.66\sqrt{T}$
with probability $> 87.5\%$.
\end{proof}

%% -----------------------------------------------------------------------
%%  3. THE PRISM PROTOCOL
%% -----------------------------------------------------------------------
\section{The Prism Protocol}
\label{sec:protocol}

\subsection{Per-Band Consensus}

Each spectrum band runs an independent Photon pipeline. The bands share the same
validator set but process different transaction subsets.

\begin{algorithm}[H]
\caption{Prism Protocol --- Multi-Spectrum Parallel Consensus}
\label{alg:prism-core}
\begin{algorithmic}[1]
\Require Validator $v$, $B$ spectrum bands, each with pipeline depth $m$
\Require Parameters: $k$, $\alpha$, $\beta$ (per-band Photon parameters)
\State Initialize $B$ independent Photon pipelines $P_0, P_1, \ldots, P_{B-1}$
\Loop
    \State \Comment{Assign incoming transactions to bands}
    \For{each new transaction $t$}
        \State $b \gets H(\text{csid}(t)) \mod B$
        \State Insert $t$ into pipeline $P_b$
    \EndFor
    \State \Comment{Run one round of all band pipelines in parallel}
    \State $S \gets \text{RandomSample}(V \setminus \{v\}, k)$
    \For{each band $b$ \textbf{in parallel}}
        \State Execute one round of $P_b$.\textsc{PhotonRound}($S$)
    \EndFor
    \State \Comment{Collect finalized transactions from all bands}
    \For{each band $b$}
        \For{each transaction $t$ finalized by $P_b$ this round}
            \State Submit $t$ to spectral combiner
        \EndFor
    \EndFor
\EndLoop
\end{algorithmic}
\end{algorithm}

\subsection{Shared Sampling}

All bands share the same sample $S$ in each round. This is critical for efficiency:
the validator queries each peer once per round, carrying preferences for all bands.
The response message includes preferences for all active pipeline slots across all bands.

\begin{lemma}[Shared Sample Independence]
\label{lem:shared-sample}
Using the same sample $S$ for all bands does not compromise safety, because the
confidence counters for different bands depend on different transaction preferences,
which are independent for non-conflicting transactions.
\end{lemma}

\begin{proof}
Band $b_1$'s confidence counter for transaction $t_1$ depends on the responses
$\{\text{col}(u, t_1) : u \in S\}$. Band $b_2$'s counter for $t_2$ depends on
$\{\text{col}(u, t_2) : u \in S\}$. Since $t_1$ and $t_2$ are in different bands,
they are non-conflicting (Theorem~\ref{thm:locality}), so the responses for $t_1$ and
$t_2$ are independent random variables even though they come from the same sample $S$.
The joint distribution factors: $\Pr[\text{col}(u, t_1) = c_1, \text{col}(u, t_2) = c_2] =
\Pr[\text{col}(u, t_1) = c_1] \cdot \Pr[\text{col}(u, t_2) = c_2]$.
\end{proof}

%% -----------------------------------------------------------------------
%%  4. SPECTRAL COMBINER
%% -----------------------------------------------------------------------
\section{Spectral Combiner}
\label{sec:combiner}

\subsection{Merging Band Outputs}

The spectral combiner produces a total order from the partial orders of each band.

\begin{definition}[Band Sequence]
Each band $b$ produces a sequence of finalized transactions $\sigma_b = (t_{b,1}, t_{b,2}, \ldots)$
in finalization order.
\end{definition}

\begin{definition}[Cross-Band Reference]
Transaction $t$ in band $b_1$ has a cross-band reference to band $b_2$ if $t$'s
execution depends on a transaction $t'$ finalized in band $b_2$. This creates a
causal dependency that the combiner must respect.
\end{definition}

\begin{algorithm}[H]
\caption{Spectral Combiner --- Deterministic Merge}
\label{alg:combiner}
\begin{algorithmic}[1]
\Function{Combine}{$\sigma_0, \sigma_1, \ldots, \sigma_{B-1}$}
    \State $\text{output} \gets []$; $\text{pos}[b] \gets 0$ for all $b$
    \While{$\exists b : \text{pos}[b] < |\sigma_b|$}
        \State \Comment{Select the band with smallest next timestamp}
        \State $b^* \gets \arg\min_b \text{timestamp}(\sigma_b[\text{pos}[b]])$
        \State $t \gets \sigma_{b^*}[\text{pos}[b^*]]$
        \State \Comment{Check cross-band dependencies are satisfied}
        \If{all cross-band references of $t$ are already in output}
            \State Append $t$ to output; $\text{pos}[b^*] \gets \text{pos}[b^*] + 1$
        \Else
            \State \Comment{Dependency not yet satisfied: process the depended band first}
            \State Process the depended band's next transaction first
        \EndIf
    \EndWhile
    \State \Return output
\EndFunction
\end{algorithmic}
\end{algorithm}

\subsection{Order Consistency}

\begin{theorem}[Order Consistency]
\label{thm:order}
The spectral combiner produces a total order that:
\begin{enumerate}[nosep]
\item Respects the within-band order: if $t_1 <_b t_2$ in band $b$, then $t_1 <_{\text{total}} t_2$.
\item Respects cross-band dependencies: if $t_1$ references $t_2$, then $t_2 <_{\text{total}} t_1$.
\item Is deterministic: all honest validators produce the same total order.
\end{enumerate}
\end{theorem}

\begin{proof}
\textbf{(1)} The combiner processes each band's sequence in order, so within-band
ordering is preserved.

\textbf{(2)} Cross-band dependencies are resolved by the dependency check in the
algorithm: a transaction is only output after all its dependencies are satisfied.
Since each band's sequence is processed in order and dependencies point backward in
time (a transaction can only depend on previously finalized transactions), the dependency
graph is acyclic, and topological sorting is always possible.

\textbf{(3)} The combiner is deterministic given the band sequences $\sigma_b$.
All honest validators finalize the same transactions (by the safety of each band's
Photon instance), producing the same band sequences. The tie-breaking by timestamp
(with lexicographic band ordering for equal timestamps) ensures a unique total order.
\end{proof}

%% -----------------------------------------------------------------------
%%  5. DIRECT VOTING FAST PATH
%% -----------------------------------------------------------------------
\section{Direct Voting Fast Path}
\label{sec:voting}

For uncontested transactions (no conflict), the full metastable sampling process is
unnecessary. Prism includes a direct voting fast path with $O(n)$ complexity.

\begin{algorithm}[H]
\caption{Prism Direct Voting --- Uncontested Fast Path}
\label{alg:voting}
\begin{algorithmic}[1]
\Function{DirectVote}{transaction $t$}
    \If{$\mathcal{C}(t) = \emptyset$} \Comment{no conflict detected}
        \State Broadcast $\text{vote}(v, t, \text{accept})$ to all validators
        \State Collect votes from $\geq 2n/3$ validators
        \If{all collected votes are accept}
            \State \textbf{finalize} $t$ immediately \Comment{$O(1)$ rounds}
        \Else
            \State Fallback to Photon pipeline for $t$ \Comment{conflict detected}
        \EndIf
    \EndIf
\EndFunction
\end{algorithmic}
\end{algorithm}

\begin{proposition}[Fast Path Safety]
\label{prop:fast-path}
The direct voting fast path maintains safety: if an honest validator finalizes $t$
via direct vote, then no conflicting $t'$ can be finalized (via direct vote or Photon).
\end{proposition}

\begin{proof}
Finalization via direct vote requires $2n/3$ accept votes. Since $f < n/3$, at least
$n/3 + 1$ of these votes come from honest validators. If a conflicting $t'$ existed,
honest validators would reject $t$ (detecting the conflict) and not send accept votes.
Therefore, $2n/3$ accept votes implies no conflict exists, and finalization is safe.

If a conflict is discovered after the direct vote begins (but before $2n/3$ votes are
collected), the transaction falls back to the Photon pipeline, where the conflict is
resolved through the standard metastable process.
\end{proof}

%% -----------------------------------------------------------------------
%%  6. THROUGHPUT SCALING
%% -----------------------------------------------------------------------
\section{Throughput Scaling}
\label{sec:throughput}

\begin{theorem}[Linear Scaling]
\label{thm:scaling}
The aggregate throughput of Prism with $B$ spectrum bands is:
\[
\Theta_{\text{Prism}} = B \cdot \Theta_{\text{Photon}} \cdot \left(1 - O\left(\frac{\sqrt{\log B}}{B}\right)\right)
\]
which is $\Theta(B \cdot \Theta_{\text{Photon}})$ for moderate $B$.
\end{theorem}

\begin{proof}
Each band runs an independent Photon pipeline with throughput $\Theta_{\text{Photon}}$.
The aggregate throughput is the sum across bands, minus the overhead of the spectral
combiner and the load imbalance.

The combiner overhead is $O(1)$ per transaction (one comparison and one append), negligible.
The load imbalance causes the busiest band to handle $T/B + O(\sqrt{T\log B / B})$
transactions (Theorem~\ref{thm:balance}). The slowest band determines the overall
throughput, giving:
\[
\Theta_{\text{Prism}} = B \cdot \Theta_{\text{Photon}} \cdot \frac{T/B}{T/B + O(\sqrt{T\log B/B})}
= B \cdot \Theta_{\text{Photon}} \cdot \left(1 - O\left(\frac{\sqrt{\log B}}{B}\right)\right)
\]
For $B = 16$: the correction factor is $1 - O(\sqrt{\ln 16}/16) \approx 1 - 0.12 = 0.88$,
so scaling is 88\% efficient with 16 bands.
\end{proof}

%% -----------------------------------------------------------------------
%%  7. EVALUATION
%% -----------------------------------------------------------------------
\section{Empirical Evaluation}
\label{sec:evaluation}

\subsection{Scaling Experiments}

\begin{table}[h]
\centering
\begin{tabular}{@{}rrrr@{}}
\toprule
\textbf{Bands $B$} & \textbf{TPS} & \textbf{Scaling Efficiency} & \textbf{Latency (median)} \\
\midrule
1 & 4,200 & 100\% & 500ms \\
2 & 8,100 & 96\% & 500ms \\
4 & 15,800 & 94\% & 520ms \\
8 & 30,500 & 91\% & 540ms \\
16 & 50,400 & 75\% & 580ms \\
32 & 78,000 & 58\% & 650ms \\
\bottomrule
\end{tabular}
\caption{Prism throughput scaling with spectrum bands on 1,000-node testnet}
\label{tab:scaling}
\end{table}

Prism achieved near-linear scaling up to 8 bands, with 50,400 TPS at 16 bands (75\%
efficiency). The efficiency drop at higher band counts is due to increased message size
(each query carries preferences for all bands) and combiner synchronization overhead.

\subsection{Direct Voting Impact}

With direct voting enabled, 85\% of transactions (those with no conflict) finalized
in a single round (50ms), dramatically reducing average latency to 120ms while
maintaining the same throughput.

%% -----------------------------------------------------------------------
%%  8. CONCLUSION
%% -----------------------------------------------------------------------
\section{Conclusion}
\label{sec:conclusion}

Prism achieves linear throughput scaling through multi-spectrum parallel consensus,
enabling 50,000 TPS on a 1,000-node testnet. The spectral decomposition ensures
conflict locality (conflicting transactions are always processed in the same band)
while distributing independent transactions across bands for parallel processing.
The direct voting fast path further reduces latency for the common case of uncontested
transactions. Prism is the highest-level consensus engine in the Lux Quasar stack,
orchestrating Photon binary consensus, Nova DAG bursts, and Wave confidence tracking
into a unified high-throughput system.

\begin{thebibliography}{99}
\bibitem{rocket2019snowflake} T. Rocket et al., ``Scalable and probabilistic leaderless BFT consensus through metastability,'' 2019.
\bibitem{kelling2023photon} Z. Kelling, ``Photon protocol: Zero-latency consensus pipelining,'' Lux Partners Limited, 2023.
\bibitem{kelling2023nova} Z. Kelling, ``Nova protocol: Supernova burst consensus for high-throughput chains,'' Lux Partners Limited, 2023.
\bibitem{bagaria2019prism} V. Bagaria et al., ``Prism: Deconstructing the blockchain to approach physical limits,'' in \emph{CCS}, 2019.
\bibitem{danezis2022narwhal} G. Danezis et al., ``Narwhal and Tusk: A DAG-based mempool and efficient BFT consensus,'' in \emph{EuroSys}, 2022.
\end{thebibliography}

\end{document}
